% ===================================================================
%  LIMINAIRES — feuillets 1 a 8 : couverture, faux-titre, page de
%  titre, verso (depot legal), AVERTISSEMENT (E. Rochelle).
%  Aucune de ces pages ne porte de folio imprime, SAUF les feuillets 6,
%  7 et 8 (suite de l'AVERTISSEMENT), qui portent respectivement
%  « — 4 — », « — 5 — » et « — 6 — » : ce sont des pages de
%  continuation, non des pages d'ouverture (meme convention que dans
%  les tableaux : voir texto/io/10-tabelo-01.tex, feuillet 7 sans folio
%  / feuillet 8 avec folio 6).
%
%  Les cles « %%K » NE SONT PAS les memes que dans texto/io/00-liminari
%  au-dela de la structure generale : la PREFACO de Guignon (ido) et
%  l'AVERTISSEMENT de Rochelle (francais) sont deux textes distincts,
%  non l'un la traduction de l'autre. La numerotation des alineas
%  releve donc l'ordre reel du francais, sans chercher a faire
%  coincider les numeros avec l'ido.
%
%  ATTENTION — LE FAC-SIMILE FRANCAIS EST UNE PRISE DE VUE, contraste
%  inegal d'une page a l'autre. Corps et interlettrages ci-dessous sont
%  ESTIMES A L'ŒIL, PROVISOIRES comme le reste du releve d'apparat ;
%  les lignes « LIVRET EXPLICATIF » / « DES » reprennent les valeurs
%  deja mesurees pour ce meme texte au feuillet 9 (texto/fr/10-tableau-
%  01.tex), le mot etant identique.
%
%  Feuillet 2 (« Sujet de chaque tableau ») : deux fragments de deux
%  lignes de la liste (tableaux 3 et 4, premiere serie) sont perdus sur
%  le fac-simile — non pas noircis, mais simplement absents d'encre
%  entre « la Fête » et « lique) » d'une part, entre « le Repas » et
%  « oce), » d'autre part. Conformement a la consigne, ces mots ne sont
%  PAS devines : ils sont marques \textit{[illisible]}. Le tableau a
%  crochets accolant deux sujets sous un meme numero n'a pas de macro
%  dediee ; il est rendu ici par une simple repetition du numero entre
%  crochets, provisoire comme le reste de l'apparat.
% ===================================================================

% --- Feuillet 1 (couverture, non foliotee) --------------------------
%  Pas d'encadrement de page. Le filet sous « E. ROCHELLE » est un
%  trait ONDULE (dents de scie fines), non un filet droit : rendu ici
%  par \VUfilet, faute de macro dediee. Pas d'ornement, pas de double
%  filet, pas de millesime sur la couverture (a la difference du
%  feuillet 3).
\begin{VUpage}[1]{}
%%K t00-apar-1 apar
\VUsaut{5.0mm}
\VUtitre{14.0pt}{60}{E. ROCHELLE}
\VUsaut{2.0mm}
\VUfilet{78mm}
\VUsaut{10.0mm}
\VUtitre{15.5pt}{50}{LIVRET EXPLICATIF}
\VUsaut{3.0mm}
\VUcentre{8.4pt}{160}{DES}
\VUsaut{4.0mm}
\VUtitre{19.5pt}{10}{Tableaux Auxiliaires Delmas}
\VUsaut{5.0mm}
\VUcentre{8.0pt}{100}{POUR}
\VUsaut{3.0mm}
\VUcentre{9.5pt}{30}{L'ENSEIGNEMENT PRATIQUE DES LANGUES VIVANTES}
\VUsaut{4.0mm}
\VUcentre{10.0pt}{0}{par l'Image et la Méthode directe}
\VUsaut{22.0mm}
\VUtitre{10.0pt}{40}{BORDEAUX}
\VUsaut{2.5mm}
\VUtitre{9.5pt}{10}{G. DELMAS, Imprimeur-Éditeur}
\VUsaut{2.0mm}
\VUcentre{8.5pt}{30}{10, RUE SAINT-CHRISTOLY}
\end{VUpage}

% --- Feuillet 2 (« Sujet de chaque tableau », non foliote) ----------
\begin{VUpage}[2]{}
%%K t00-apar-2 apar
\VUsaut{3.0mm}
\VUtitre{18.0pt}{20}{Tableaux Auxiliaires Delmas}
\VUsaut{2.5mm}
\VUtitre{10.5pt}{40}{EN SIX LANGUES}
\VUsaut{4.0mm}
\VUcentre{9.5pt}{0}{\textit{(Anglais, Allemand, Italien, Espagnol, Russe, Français)}}
\VUsaut{4.0mm}
\VUornamento{9pt}
\VUsaut{4.5mm}
\VUtitre{13.0pt}{40}{SUJET DE CHAQUE TABLEAU}
\VUsaut{4.0mm}
\VUtitre{10.5pt}{20}{PREMIÈRE SÉRIE}
\VUsaut{2.0mm}
\noindent\textsc{Tableau}\nl
N\textsuperscript{o}\ \ 1 --- L'École, la Classe, les Nombres.\nl
2 --- La Récréation, les Jeux, le Corps humain.\nl
[3] L'Enfance (le Baptême).\nl
[3] La Jeunesse (la Fête \textit{[illisible]} lique).\nl
[4] L'Age mûr (le Repas \textit{[illisible]} oce),\nl
[4] La Vieillesse (la Visite \textit{[illisible]} grand-père).\nl
5 --- La Maison extérieure, les Ouvriers, les Outils.\nl
6 --- La Maison intérieure, les Meubles, les Ustensiles.
\VUsaut{4.0mm}
\VUtitre{10.5pt}{20}{DEUXIÈME SÉRIE}
\VUsaut{2.0mm}
\noindent N\textsuperscript{o}\ \ [7] Le Village (en Hiver).\nl
[7] La Maison rustique (au Printemps).\nl
[8] La Moisson (en Été).\nl
[8] Les Vendanges, la Pêche (en Automne).\nl
[9] La Montagne et la Station thermale.\nl
[9] La Forêt, la Chasse.\nl
10 --- La Mer, la Plage, le Port.
\VUsaut{4.0mm}
\VUtitre{10.5pt}{20}{TROISIÈME SÉRIE}
\VUsaut{2.0mm}
\noindent N\textsuperscript{o}\ 11 --- La Ville et ses Monuments (un Incendie).\nl
12 --- La Gare, les Voyages, les Bateaux.\nl
13 --- L'Hôtel, le Restaurant, le Café.\nl
14 --- La Rue, les Commerçants, les Métiers.\nl
15 --- Le Marché aux comestibles (la Nourriture).\nl
16 --- Un grand Magasin (Objets de toutes sortes).
\end{VUpage}

% --- Feuillet 3 (page de titre, non foliotee) -----------------------
%  Meme composition que le feuillet 1, au mot pres : l'ordre des deux
%  membres du sous-titre est INVERSE (« Par la Méthode directe et par
%  l'Image », majuscule initiale), un ornement (volute a crossettes,
%  non reproduit) suit le sous-titre, et le bloc d'adresse porte un
%  millesime que la couverture n'a pas.
\begin{VUpage}[3]{}
%%K t00-apar-3 apar
\VUsaut{5.0mm}
\VUtitre{14.0pt}{60}{E. ROCHELLE}
\VUsaut{2.0mm}
\VUfilet{78mm}
\VUsaut{10.0mm}
\VUtitre{15.5pt}{50}{LIVRET EXPLICATIF}
\VUsaut{3.0mm}
\VUcentre{8.4pt}{160}{DES}
\VUsaut{4.0mm}
\VUtitre{19.5pt}{10}{Tableaux Auxiliaires Delmas}
\VUsaut{5.0mm}
\VUcentre{8.0pt}{100}{POUR}
\VUsaut{3.0mm}
\VUcentre{9.5pt}{30}{L'ENSEIGNEMENT PRATIQUE DES LANGUES VIVANTES}
\VUsaut{4.0mm}
\VUtitre{10.0pt}{10}{Par la Méthode directe et par l'Image}
\VUsaut{10.0mm}
\VUornamento{9pt}
\VUsaut{14.0mm}
\VUtitre{10.0pt}{40}{BORDEAUX}
\VUsaut{2.5mm}
\VUtitre{9.5pt}{10}{G. DELMAS, Imprimeur-Éditeur.}
\VUsaut{2.0mm}
\VUcentre{8.5pt}{30}{10, RUE SAINT-CHRISTOLY, 10.}
\VUsaut{4.0mm}
\VUfilet{8mm}
\VUsaut{3.0mm}
\VUtitre{9.5pt}{20}{1903}
\end{VUpage}

% --- Feuillet 4 (verso de la page de titre, non foliote) ------------
%  Grande page blanche : seuls la mention de depot legal (italique) et
%  le nom de l'editeur, signe, apparaissent en bas de page. La
%  signature manuscrite qui suit « L'Éditeur, » sur cet exemplaire
%  n'est pas de la matiere imprimee : non transcrite, comme la marque
%  d'imprimeur du feuillet 3 ido.
\begin{VUpage}[4]{}
%%K t00-apar-4 apar
\VUsaut{140.0mm}
{\centering\itshape\fontsize{9.0pt}{11.0pt}\selectfont
Le dépôt du présent ouvrage a été effectué confor\cc
mément aux lois françaises et aux conventions\nl
internationales.\par
Tous droits réservés. --- Reproduction interdite.\par}
\VUsaut{6.0mm}
{\raggedleft\fontsize{10.0pt}{12.0pt}\selectfont L'Éditeur,\par}
\end{VUpage}

% --- Feuillet 5 (ouverture de l'AVERTISSEMENT, non foliote) ---------
\begin{VUpage}[5]{}
%%K t00-tit-1 sub
\VUsaut{16.0mm}
\VUtitre{19.0pt}{40}{Avertissement}
\VUsaut{4.0mm}
\VUfilet{12mm}
\VUsaut{8.0mm}

%%K t00-01-1 p
\VUblancAlinea
Beaucoup de nos collègues, qui ont adopté\nl
dès le premier jour nos \VUgras{Tableaux auxiliai}\cc
\VUgras{res}, nous ont demandé de leur rendre la tâche\nl
plus facile en expliquant \VUgras{ce que nous avions}\nl
\VUgras{voulu mettre dans chaque tableau.}

%%K t00-02-1 p
\VUblancAlinea
C'est à leur intention que nous avons rédigé\nl
ce petit livret, qui donne une forme vivante au\nl
vocabulaire qui accompagne le dessin.

%%K t00-03-1 p
\VUblancAlinea
Il permettra au professeur de voir d'un seul\nl
coup d'œil tout ce que contient chaque tableau\nl
et quel est le parti qu'il peut en tirer dans\nl
l'application de la \VUgras{méthode directe} telle\nl
qu'elle a été définie pour la France d'une façon\nl
désormais officielle dans les programmes et\nl
commentée par MM. les Inspecteurs généraux\nl
des langues vivantes dans leurs conférences\nl
pédagogiques.

%%K t00-04-1 p
\VUblancAlinea
\VUgras{Les Tableaux auxiliaires Delmas} ne\nl
constituent pas une \VUgras{méthode}. Ils n'ont qu'une\nl
prétention, c'est d'être un auxiliaire commode,\nl
un simple instrument, qui permet au professeur\nl
de donner, en 6\textsuperscript{e} et en 5\textsuperscript{e}, un enseignement tout
\parplein
\end{VUpage}

% --- Feuillet 6 (suite de l'AVERTISSEMENT, folio 4) ------------------
\begin{VUpage}[6]{4}
%%K t00-04-1 p suite
\VUcontinue
oral, en lui fournissant d'innombrables \VUgras{sujets}\nl
\VUgras{de conversation} sur le vocabulaire prescrit\nl
pour ces deux classes.

%%K t00-05-1 p
\VUblancAlinea
Il ne faut pas oublier, en effet, que \VUgras{la vraie}\nl
\VUgras{méthode directe est essentiellement}\nl
\VUgras{personnelle,} et nous nous garderions bien\nl
d'empiéter en quoi que ce soit sur cette \VUgras{initia}\cc
\VUgras{tive} féconde qui est si généreusement laissée\nl
au professeur. Ce n'est pas à nous de lui fixer\nl
ce qu'il doit choisir dans le vocabulaire pres\cc
crit. Il reste entièrement libre de son choix\nl
comme aussi de l'ordre dans lequel il fera\nl
étudier les différents tableaux. Nous avons suivi\nl
la classification officielle, qui n'est impérative\nl
pour personne.

%%K t00-06-1 p
\VUblancAlinea
Ceci dit, voyons comment nous allons nous\nl
servir du tableau.

%%K t00-07-1 p
\VUblancAlinea
Le professeur, après l'avoir étudié lui-même\nl
à l'aide du livret ci-joint, commentera en\nl
phrases très simples, en ne se servant, bien\nl
entendu, que de la langue qu'il enseigne, la\nl
scène que représente le dessin. Après l'avoir\nl
décrite dans ses grandes lignes, il entrera dans\nl
plus de détails et la divisera en petites scènes\nl
secondaires. Il s'assurera par d'habiles ques\cc
tions qu'il a été bien compris de tous, et que\nl
la plupart des élèves sont capables de répéter\nl
ce qu'il vient de dire.

%%K t00-08-1 p
\VUblancAlinea
Pour fixer ces notions acquises d'une façon\nl
tout orale, il y aura plusieurs procédés. Celui\nl
que semblent recommander les instructions
\parplein
\end{VUpage}

% --- Feuillet 7 (suite de l'AVERTISSEMENT, folio 5) ------------------
\begin{VUpage}[7]{5}
%%K t00-08-1 p suite
\VUcontinue
officielles consistera à écrire au tableau les\nl
phrases que l'on veut faire retenir. Les élèves\nl
les recopieront sur leur cahier. Malheureuse\cc
ment l'expérience nous apprend que, malgré\nl
tous les efforts du professeur pour revoir soi\cc
gneusement les cahiers, de nombreuses erreurs\nl
se seront glissées dans ce texte, erreurs qui ris\cc
queront de se fixer dans la mémoire de l'enfant.

%%K t00-09-1 p
\VUblancAlinea
C'est alors qu'interviendra le vocabulaire an\cc
nexé au tableau.

%%K t00-10-1 p
\VUblancAlinea
Ce vocabulaire, \VUgras{préalablement séparé de}\nl
\VUgras{la gravure, n'a pas été apporté en classe.}\nl
La vue du mot écrit n'a donc pas pu troubler\nl
l'image auditive qui s'est formée dans l'esprit\nl
de l'élève. Mais l'enfant le retrouve chez lui, ou\nl
en étude, ce précieux \VUgras{memento} qui, à l'aide\nl
de ses numéros, va lui permettre de rechercher\nl
l'orthographe exacte et le sens du mot mal\nl
compris, \VUgras{sans passer par l'intermédiaire}\nl
\VUgras{de la langue maternelle.}

%%K t00-11-1 p
\VUblancAlinea
Le Tableau auxiliaire présente donc un double\nl
avantage. Il rend au professeur la tâche moins\nl
pénible. Il fait, pour l'élève, de l'étude de la\nl
langue étrangère un travail amusant.

%%K t00-12-1 p
\VUblancAlinea
Il offre encore une utilité incontestable, c'est\nl
de fournir d'innombrables sujets de \VUgras{devoirs}\nl
\VUgras{faciles} : descriptions, narrations, lettres, etc.

%%K t00-13-1 p
\VUblancAlinea
{\itshape Par exemple~: (Tableau n\textsuperscript{o} 1, allemand)~:\nl
Décrire une partie de la salle de classe en\nl
employant les prépositions auf, in, neben,\nl
vor et hinter.\par}
\end{VUpage}

% --- Feuillet 8 (fin de l'AVERTISSEMENT, folio 6) ---------------------
\begin{VUpage}[8]{6}
%%K t00-14-1 p
\VUblancAlinea
Tels sont, à notre avis, les services que peu\cc
vent rendre à nos collègues et à leurs élèves\nl
les \VUgras{Tableaux auxiliaires Delmas.}

%%K t00-15-1 p
\VUblancAlinea
Comme toutes les choses humaines, ils sont\nl
essentiellement perfectibles, et leur éditeur est\nl
décidé à ne reculer devant aucun sacrifice pour\nl
les rendre plus clairs et plus pratiques.

%%K t00-16-1 p
\VUblancAlinea
C'est pourquoi nous prions nos collègues qui\nl
s'en servent de ne pas hésiter à nous signaler\nl
toutes les améliorations qui leur paraîtront\nl
nécessaires.
\VUsaut{6.0mm}
{\raggedleft\fontsize{11.0pt}{13.0pt}\selectfont\textsc{E. Rochelle,}\par}
{\raggedleft\itshape\fontsize{9.5pt}{11.5pt}\selectfont Professeur au Lycée de Bordeaux.\par}

\VUsaut{18.0mm}
%%K t00-noto-1 noto
{\fontsize{9.0pt}{11.0pt}\selectfont
\noindent\textsc{Nota.} --- Nous avons imprimé en \VUgras{caractères gras} les\nl
\textit{substantifs} qui se trouvent dans le vocabulaire des\nl
Tableaux en les faisant suivre de leur numéro. Mais nous\nl
engageons le professeur à attirer plus spécialement\nl
l'attention de l'élève sur les nombreux \textit{verbes} introduits\nl
dans la description. Les verbes constituent, en effet, la\nl
partie la plus importante de la langue, celle qui présente\nl
le plus de difficultés.\par}
\end{VUpage}
